\documentclass[11pt,a4paper]{article}
\usepackage[margin=0.85in,top=1in,bottom=1in]{geometry}
\usepackage{amsmath,amssymb,amsfonts}
\usepackage{graphicx}
\usepackage{float}
\usepackage{enumitem}
\usepackage{microtype}
\usepackage{caption}
\usepackage[hidelinks]{hyperref}
\usepackage[most,breakable]{tcolorbox}
\tcbuselibrary{breakable,skins}
\usepackage{fontspec}
\setmainfont{DejaVu Serif}
\setsansfont{DejaVu Sans}
\setmonofont{DejaVu Sans Mono}
\captionsetup{font=small,labelfont=bf,skip=4pt}
\setlist{leftmargin=1.5em,topsep=2pt}

%% Breakable card — prevents content cutoff across pages
\newtcolorbox{solcard}[3]{
  breakable, enhanced,
  colback=white, colframe=gray!50, arc=2pt,
  title={\textbf{#1}\hfill\textsf{\small #3\,pts}},
  fonttitle=\normalsize\sffamily\bfseries,
  before upper={\small\textit{#2}\par\smallskip\hrule height 0.4pt\smallskip},
  top=4pt, bottom=6pt, left=7pt, right=7pt,
  parbox=false,
}

\title{\textbf{The main thing is not to oversalt --- Solution}}
\author{\normalsize W21 (2021) — English translation}
\date{}
\begin{document}\maketitle\vspace{-0.5em}
\begin{solcard}{A1}{Measure the resistance of the limiting resistor $R_1$ with an ohmmeter (see equipment). The approximate value of 10 ohms is indicated below, but use the more accurate value you obtained in your calculations.}{??}

Let's measure the resistance of the limiting resistor $R_1 = 10,1~\mathrm{\Omega}$ with a multimeter in ohmmeter mode.

\end{solcard}

\begin{solcard}{A2}{Experimentally confirm the validity of the proposed model. To do this, study the current-voltage characteristics of the solution at 5 different distances between the electrodes.}{7.50}

We assemble a measuring circuit consisting of a battery connected in series, a control unit and a syringe with a salt solution. We connect a voltmeter to the electrodes of the syringe. We connect the second voltmeter in parallel with the reference resistance $R_1$ to measure the current in the circuit.


The entire set of processes on the cathode and anode is proposed to be represented in the form of a series circuit consisting of an external voltage applied to it $U_0$, resistance of the electrolyte volume $R$, contact EMF $U_к$ and contact resistance $r_к$ (Fig. 1). It is assumed that the value $R$ is directly proportional to the resistivity of the electrolyte $\rho$, the distance between the electrodes $L$ and inversely proportional to the cross-sectional area of ​​the electrolytic bath $S$.


The theoretical current-voltage characteristic of the electrolytic bath within the framework of this model has the form $\rho$ and $r_k$.


We take the current-voltage characteristic of the solution. The tables show the results of measuring the current-voltage characteristics of the solution at 5 different volumes of liquid in the syringe $V$ (and the corresponding distances $L$ between the electrodes).


$U_{10}$ , mV $U$, mV $I$, mA $V$, ml $L$, cm 290 1740 28.7 4 1.3 320 1860 31.7 450 2060 44.6 650 2460 64.4 770 2670 76.2 990 3080 98.0 1080 3420 106.9


$U_{10}$ , mV $U$, mV $I$, mA $V$, ml $L$, cm 870 3750 86.1 8 2.6 800 3570 79.2 740 3340 73.3 540 2920 53.5 392 2300 38.8 294 2040 29.1 269 1970 26.6


$U_{10}$ , mV $U$, mV $I$, mA $V$, ml $L$, cm 254 2160 25.1 12 3.9 269 2230 26.6 330 2450 32.7 490 3100 48.5 610 3550 60.4 670 3790 66.3 730 3950 72.3


$U_{10}$ , mV $U$, mV $I$, mA $V$, ml $L$, cm 574 4140 56.8 16 5.2 578 4010 57.2 520 3800 51.5 475 3400 47.0 425 3180 42.1 319 2650 31.6 266 2400 26.3


$U_{10}$ , mV $U$, mV $I$, mA $V$, ml $L$, cm 213 2600 21.1 20 6.5 289 3070 28.6 360 3410 35.6 451 3950 44.7 501 4150 49.6 523 4250 51.8


We build the current-voltage characteristics for all experiments

\begin{figure}[H]
  \centering
  \includegraphics[width=0.72\textwidth]{E2_S_EN_assets/fig1.jpg}
\end{figure}
\begin{figure}[H]
  \centering
  \includegraphics[width=0.72\textwidth]{E2_S_EN_assets/fig2.jpg}
\end{figure}
\end{solcard}

\begin{solcard}{A3}{Estimate the magnitude of the contact potential difference $U_\text{to}$.}{1.00}

Using the graph, we determine: 1. From the point of intersection of the graphs and the voltage axis, the value of the contact potential difference $U_к = 1,2~\mathrm{V}$. 2. From the angular coefficient (see formula $(2)$) the resistance of the electrolyte volume together with the resistance of the contact areas $R_{tot}$ for 5 values ​​$L$.

\end{solcard}

\begin{solcard}{A4}{Determine the value of contact resistance $r_\text{to}$.}{2.50}

The results of processing experimental data for five values ​​of $L$ are presented in the table and figure.


No. $L$, m $1/R_{tot}$ , 1/Ohm $R_{tot}$ , Ohm 1 0.065 0.016 62.5 2 0.052 0.020 50.0 3 0.039 0.026 38.5 4 0.026 0.033 30.3 5 0.013 0.049 20.4


From the graph we find $r_к = R_{tot}(0) = 9,2~\mathrm{\Omega}$.

\begin{figure}[H]
  \centering
  \includegraphics[width=0.72\textwidth]{E2_S_EN_assets/fig3.jpg}
\end{figure}
\end{solcard}

\begin{solcard}{A5}{Calculate the specific conductivity $\sigma$ of a three percent solution of sodium chloride in water.}{1.00}

Cross-sectional area of ​​the piston $S = 3,1·10^{-4}~\mathrm{m}^2$ (can be found by measuring the diameter of the piston or the ratio of volume to length). Within the framework of the used model $ ΔR_{tot}/ΔL = ρ/S$. From the slope of the straight line $R_{tot}(L)$ we find the resistivity of the salt solution used $ρ = 0,24~ \mathrm{\Omega}·м$. That is, $σ = 1/ρ =  4,2 ~(Ом·м)^{-1}$.

\end{solcard}

\begin{solcard}{A6}{Calculate the mobility $\mu$ of sodium and chlorine ions in this solution.}{3.00}

Specific conductivity of the electrolyte $σ = 1/ρ = q_1n_1µ_1 + q_2n_2µ_2$ where $q_1$ and $q_2$ are the modules of the electrical charges of positive and negative ions, respectively (in the case of sodium and chlorine ions $q_1 =  q_2 = е$, where $е$ is the elementary charge), $n_1$ and $n_2$ are their concentrations (in our case they are the same and equal to $n$), $µ_1$ and $µ_2$ are the mobilities of negative and positive ions in solution, which, according to the conditions of the problem, are also equal to each other $µ_1 = µ_2  = µ$.  Finally, $$µ=\frac{σ}{2еn}$$


We calculate the concentration of sodium (or chlorine) ions in a three percent solution of table salt $NaCl$ $$n =N/V=\frac{m_{NaCl}N_a\rho_{раств}}{M_{NaCl}m_{раств}}=0,03\frac{N_a\rho_{раств}}{M_{NaCl}}= 3,2·10^{26} ~\mathrm{m}^{-3}$$.


We calculate the mobility: $$µ = \frac{σ}{(2еn)} = 4,0·10^{-8}\frac{м^2}{В·с}$$, which is in good agreement with the literature data.

\end{solcard}

\end{document}